\documentclass[11pt]{article}
\usepackage[margin=1in]{geometry}
\usepackage{amsmath,amssymb,amsfonts,mathtools,bm}
\usepackage[T1]{fontenc}
\usepackage[utf8]{inputenc}
\title{Inline and Display Mathematics}
\author{Source-backed grouped formula sample}
\date{}
\begin{document}
\maketitle
\section{Expressions}
\paragraph{Expression 1.} The following expression is taken from a source-backed formula corpus.

\begin{align*}\tilde{\hat{\Omega}}_{\underline{x}}{}^{(0)\hat{a}\hat{b}} =-\hat{\Omega}_{\underline{x}}{}^{(0)\hat{a}\hat{b}} /\hat{G}_{\underline{xx}}^{(0)}\, ,\hspace{.5cm}\tilde{\hat{\Omega}}_{\mu}{}^{(0)\hat{a}\hat{b}} =\hat{\Omega}_{\mu}{}^{(0)\hat{a}\hat{b}}-\hat{\Omega}_{\underline{x}}{}^{(0)\hat{a}\hat{b}}\hat{G}_{\underline{x}\mu}^{(0)}/\hat{G}^{(0)}_{\underline{xx}}\, ,\end{align*}

\paragraph{Expression 2.} The following expression is taken from a source-backed formula corpus.

\begin{align*}\left[ \bar \gamma ^2\bar \gamma ^0\bar \gamma ^1\partial_\vartheta +\frac{\bar \gamma ^3\bar \gamma ^0\bar \gamma ^1}{\sin \vartheta}\partial _\varphi +\bar \gamma ^2\bar \gamma ^0\bar \gamma ^1S\partial_\vartheta S^{-1}+\frac{\bar \gamma ^3\bar \gamma ^0\bar \gamma ^1}{\sin\vartheta }S\partial _\varphi S^{-1}\right] \tilde \Phi =-i\kappa \tilde\Phi\end{align*}

\paragraph{Expression 3.} The following expression is taken from a source-backed formula corpus.

\begin{align*}\frac{1}{\sqrt{-g^{\prime}}}\partial_{\mu}(\sqrt{-g^{\prime}}g^{\prime\mu\nu}\partial_{\nu}\phi)+ \frac{1}{\chi}\left[\frac{dV_{1}}{d\phi}-\frac{1}{\chi}\frac{dV_{2}}{d\phi}+\frac{\alpha}{M_{p}}P(\varphi)e^{\alpha\phi /M_{p}}\right]=-\frac{h\gamma}{\sqrt{2}M_{p}}\overline{\Psi}^{\prime}\Psi^{\prime} \varphi \frac{e^{\gamma\phi/M_{p}}}{\chi^{3/2}}; \end{align*}

\paragraph{Expression 4.} The following expression is taken from a source-backed formula corpus.

\begin{align*}\left| \sum_{\alpha \beta \gamma \delta}\langle \Phi_{\alpha} \Phi_{\beta} \Phi_{\gamma} \Phi_{\delta} \rangle_c\left. \right/ \sum_{\alpha \beta \gamma \delta}\langle \Phi_{\alpha} \Phi_{\beta} \Phi_{\gamma} \Phi_{\delta} \rangle\right|=\left| \langle \Phi^4(x) \rangle_c \left. \right/\langle \Phi^4(x) \rangle \right|\; (n=4) \; ,\end{align*}

\paragraph{Expression 5.} The following expression is taken from a source-backed formula corpus.

\begin{align*}{\bar{{\cal G}}}_{({\rm I})}^{\left( \rho \right) a}\Gamma ^{\left( \rho\right) }=\Delta _{{\rm cl}({\rm I})}^{\left( \rho \right) a}+O\left( \rho^2\right) \,\,,\;\;\;{\bar{{\cal G}}}_{({\rm II})}^{\left( \rho \right)a}\Gamma ^{\left( \rho \right) }=\Delta _{{\rm cl}({\rm II})}^{\left( \rho\right) a}+O\left( \rho ^2\right) \,\,, \end{align*}
\end{document}
